\section{Glossary}
\label{sec:glossary}

Definitions as used \emph{in this project} (handout
\texttt{docs/CDP\_explainer\_students.pdf} and Loftus, Bynum \& Hansen,
arXiv:2303.04209). Each row points to the section where the term is
developed.

\begin{longtable}{@{}p{3.1cm}p{7.4cm}p{2.4cm}@{}}
\toprule
\textbf{Term} & \textbf{Means, here} & \textbf{See} \\
\midrule
\endhead
ECM & ``Explanatory causal model'' --- the analyst's DAG over the model's
  \emph{features} (not a claim about the real world, just the assumption
  used to define the explanation) & \S\ref{sec:bigpicture} \\
$Y^*$, constructed outcome & The fitted model's output, $Y^*=f(X)$, treated
  as if it were an outcome variable for causal-inference purposes &
  \S\ref{sec:bigpicture} \\
PDP & Partial Dependence Plot (Friedman 2001): average $f$ over the data
  with feature $A$ set to $a$, other features left at recorded values &
  \S\ref{sec:bigpicture} \\
ICE & Individual Conditional Expectation plot: one PDP-style curve per
  observation, not averaged & \S\ref{sec:papers} \\
PDP-based feature importance & Variance (numerical features) or range$/4$
  (categorical) of the PDP curve's own values --- how much the curve moves,
  as a single number; blind to interaction effects & \S\ref{sec:bigpicture} \\
TDP & Total Dependence Plot: $\operatorname{do}(A=a)$, let downstream
  features respond & \S\ref{sec:bigpicture} \\
NDDP & Natural Direct Dependence Plot: set $A=a$, hold downstream features
  at natural values --- provably equals the PDP & \S\ref{sec:bigpicture} \\
NIDP & Natural Indirect Dependence Plot: keep your own $A$, but downstream
  features move to where they'd be under $A=a$ & \S\ref{sec:bigpicture} \\
PCDP & Partially Controlled Dependence Plot: $\operatorname{do}(A=a)$
  \emph{and} fix another feature & \S\ref{sec:bigpicture} \\
Reduced form, $Q(dr\mid a,w)$ & Conditional law of the non-parent features
  $R$ given $(A,W)$ --- learnable straight from data, no ECM equations
  needed & \S\ref{sec:identification} \\
Bridge mean & Fulcher et al.'s object $\Psi(a)=E[Y(A,Z(a))]$; turns out to
  be exactly the NIDP & \S\ref{sec:identification} \\
AIPW & Augmented inverse-probability weighting: doubly-robust estimator,
  consistent if \emph{either} the propensity or the outcome regression is
  right & \S\ref{sec:estimation} \\
Cross-fitting & Fit nuisance models on one data split, evaluate on another,
  swap --- lets flexible ML be used for nuisances and still get valid Wald
  intervals & \S\ref{sec:estimation} \\
Efficient covariate set & Not just \emph{a} valid adjustment set (e.g.\
  parents $W$) but the one minimizing variance --- here, all
  non-descendants of $A$ when $f$ uses every feature &
  \S\ref{sec:estimation} \\
Non-overlap ratio & Diagnostic (Bao \& Schomaker 2026) for how much the
  curve at $a$ relies on extrapolation vs.\ observed data overlap &
  \S\ref{sec:propagation} \\
Tipping strength & Cinelli--Hazlett robustness value, reparameterized per
  ECM edge: how strong an unmeasured confounder $U$ would need to be to flip
  the curve's slope & \S\ref{sec:propagation} \\
Propagation & Alternative to adjustment: fit every ECM structural equation
  and simulate counterfactuals person-by-person (what the original CDP
  paper does) & \S\ref{sec:propagation} \\
\bottomrule
\caption{Project glossary.}
\label{tab:glossary}
\end{longtable}

\emph{Add rows as new project-specific terms are met --- don't let this go
stale.}
